\documentclass[12pt,letterpaper]{article}
\usepackage{amsmath}
\usepackage{amsthm}
\usepackage{amsfonts}
\usepackage{amssymb}
\usepackage{amscd}
\usepackage{enumerate}
\usepackage{fancyhdr}
\usepackage{mathrsfs}
\usepackage{bbm}
\usepackage{framed}
\usepackage{mdframed}
\usepackage{cancel}
\usepackage{mathtools}
\usepackage{verbatim}
\usepackage{hyperref}
\usepackage[letterpaper,voffset=-.5in,bmargin=3cm,footskip=1cm]{geometry}
\setlength{\parindent}{0.0in}
\setlength{\parskip}{0.1in}
\allowdisplaybreaks
\headheight 15pt
\headsep 10pt
\newcommand\N{\mathbb N}
\newcommand\Z{\mathbb Z}
\newcommand\R{\mathbb R}
\newcommand\Q{\mathbb Q}
\newcommand\lcm{\operatorname{lcm}}
\newcommand\setbuilder[2]{\ensuremath{\left\{#1\;\middle|\;#2\right\}}}
\newcommand\E{\operatorname{E}}
\newcommand\V{\operatorname{V}}
\newcommand\Pow{\ensuremath{\operatorname{\mathcal{P}}}}

\DeclarePairedDelimiter\ceil{\lceil}{\rceil}
\DeclarePairedDelimiter\floor{\lfloor}{\rfloor}
\newcommand\hint[1]{\textbf{Hint}: #1}
\newcommand\note[1]{\textbf{Note}: #1}
\newcounter{enum22i}
\newenvironment{22enumerate}{\begin{enumerate}[a.]\itemsep0em\setcounter{enumi}{\value{enum22i}}}{\setcounter{enum22i}{\value{enumi}}\end{enumerate}}
\newenvironment{22itemize}{\begin{itemize}\itemsep0em}{\end{itemize}}
\fancypagestyle{firstpagestyle} {
  \renewcommand{\headrulewidth}{0pt}
  \lhead{\textbf{CSCI 0220}}
  \chead{\textbf{Discrete Structures and Probability}}
  \rhead{M. Littman}
}

\fancypagestyle{fancyplain} {
  \renewcommand{\headrulewidth}{0pt}
  \lhead{\textbf{CSCI 0220}}
  \chead{Homework 10}
  \rhead{\textit{ }}
}
\pagestyle{fancyplain}
\usepackage{tikz-cd}
\usepackage{tikz}
\usepackage{graphicx}

\begin{document}
  \thispagestyle{firstpagestyle}
  \begin{center}
    {\huge \textbf{Homework 10}}

    {\large Due: \textit{Monday, August 2nd at 11:59pm EDT}}
  \end{center}


     All homeworks are due the Monday after their release at 11:59pm EDT on Gradescope.

    Please do not include any identifying information about yourself in the handin, including your Banner ID.

    Be sure to fully explain your reasoning and show all work for full credit. We recommend checking out the sample proofs on the course website to see the level of rigor for which we're looking.
  

\subsection*{Problem 1}
{
{\setcounter{enum22i}{0}

     An interplanetary spaceflight map can be represented as a graph $G = (V,E)$, where $V$ is the set of all planetary spaceports and $(u,v) \in E$ if there is a spaceflight from $u$ to $v$. Suppose Pluto simulates a random spaceflight map of the solar system where, for each pair of spaceports, there is a spaceflight available between them  (in either direction) with probability $p$.
  
    \begin{22enumerate}
      \item Let random variable $F$ be the total number of spaceflights. Find the expected value of $F$, justifying your answer.
      \item Pluto wants to go on a tour of the solar system. Let random variable $H$ be the number of Hamiltonian tours in $G$. Find the expected value of $H$, justifying your answer.
      
      \note{A tour is defined by the order of the cycle rather than the order from some arbitrary starting vertex. Hence, the tour $(a,b,c,a)$ is the same tour as $(b,c,a,b)$.}
      \item Define the random variable $T$ to be the number of triplets of spaceports $(a,b,c)$ such that there is a path of exactly 2 spaceflights from $a$ to $c$ through $b$. This means that a spaceflight exists from $a$ to $b$ and from $b$ to $c$, but not from $a$ to $c$ directly. Find the expected value of $T$, justifying your answer.\\
      \hint{Use indicator random variables!}
    \end{22enumerate}

    
}
}
\newpage
\subsection*{Problem 2}
{
{\setcounter{enum22i}{0}
Using strong induction, prove that if $G$ is a simple graph with $n$ vertices, $k$ connected components, and no cycles, then $G$ has $n-k$ edges.

\note{Remember to use \textit{build-down} induction! This induction proof requires a proof on two variables. Define the predicate $P(n,k)$ where $n\geq k$ and in the inductive step show that both $P(i+1,j)$ and $P(i,j+1)$ are true (an increment on both variable independently!).}

\hint{If the connected components have no cycles, what must they be?}
}
}


\subsection*{Problem 3}
{
\newcommand\Cube{\textsf{Hypercube}}
{\setcounter{enum22i}{0}

      Define a \emph{uniformly-random $k$-coloring}
      $f : V(G) \to \{1,2,\dots,k\}$ as a coloring of $G$ (where $|V(G)|=n$) where each vertex $v \in V(G)$ is assigned one of the $k$ colors uniformly at random. This might produce either a proper or improper coloring of $G$. A graph is \emph{properly colored} if each vertex in the graph is assigned a color such that for all edges $(u,v)$, $u$ and $v$ are assigned \emph{different} colors. We say a graph is $n$-colorable if there exists a way to \emph{properly color} a graph using $n$ colors.

      \begin{22enumerate}
        \item In terms of $n$ and $k$, how many $k$-colorings of $G$ are there? Include both proper and improper colorings and justify your answer.
        \item Let $K_n$ be the complete graph on $n$ vertices, with
          $n \geq 3$.

          \begin{enumerate}[i]
            \item For what values of $k$ do
              there exist proper $k$-colorings of $K_n$?
            \item What is the probability that a uniformly-random $k$-coloring
            $f$ is a proper $k$-coloring of $K_n$? You need only consider $k$'s such that a
            proper $k$-coloring of $K_n$ is possible.
          \end{enumerate}
      \item
      For $n \geq 1$, define the graph $\Cube_{n}$ as follows. The vertex set is $\{0,1\}^{n}$ and, for binary strings of length $n$ denoted by $u$ and $v$, $\{u, v\}$ is an edge of $\Cube_n$ if and only if $u$ and $v$ differ in exactly one position. For example, $\Cube_{1}$ is a single edge, $\Cube_{2}$ is a square, and $\Cube_{3}$ is a 3-dimensional cube! \note{If it aids your understanding, try drawing these small examples! The name may become clearer to you at that point.}

      What is the minimum number of colors $k$ that are needed to properly $k$-color $\Cube_{n}$?
      \end{22enumerate}
    
}
}
\newpage
\subsection*{Mind Bender (Extra Credit)}
{
{\setcounter{enum22i}{0}

      
        Recall the graph coloring definitions from Problem 3. We say a graph is $n$-colorable if there exists a way to \emph{properly color} a graph using $n$ colors.
      
        A graph is \emph{planar} if it can be drawn on a plane in such a way  that its edges intersect only at vertices. In other words, we can draw the graph on a piece of paper in such a way that no two edges overlap.

        The 4-color theorem says that any planar graph can be 
        properly colored using only 4 colors. This theorem is famously difficult to prove... but you can prove the 6-color theorem right now! Let $G = (V,E)$ be a simple, connected, planar graph on at least three vertices. We will be using this graph for the rest of the problem.
        \begin{22enumerate}
        \item
        Recall from lecture that $\sum_{v \in V} \deg(v) = 2|E|$. Show that the average (mean) degree of vertices in $V$ is strictly
        less than 6. 

        \hint{You may use without proof that, for any simple,
        connected, planar graph on at least 3 vertices, $|E| \leq 3|V| - 6$.}
      \item
        Show by contradiction that there exists a vertex $v \in V$ such that
         $\deg(v) \leq 5$.
      \item
        Prove, by induction on the number of vertices, that $G$ is 6-colorable.

        \note{Remember to use build-down induction!}
      \end{22enumerate}
    
}
}

\end{document}